\begin{exercice*} 
    Sur la figure ci-dessous :
    \begin{itemize}
        \item Les points $O$, $P$, $T$ sont alignés.
        \item Les points $O$, $R$, $F$ sont alignés.
    \end{itemize}
    \def\width{8}
    \def\hauteur{5}
    \begin{tikzpicture}[scale=0.9]        
        \draw[step=2mm, line width=0.1mm, black!30!white] (0,0) grid (\width,\hauteur);
        \draw[step=1cm, line width=0.3mm, black!50!white] (0,0) grid (\width,\hauteur);
        \coordinate (O) at (1,1);
        \coordinate (P) at (1,2);
        \coordinate (T) at (1,4);
        \coordinate (R) at (3,1);
        \coordinate (F) at (7,1);
        \tkzDrawSegment(O,T);
        \tkzDrawSegment(T,F);
        \tkzDrawSegment(F,O);
        \tkzDrawSegment(P,R);
        \tkzLabelPoints[above left](T);
        \tkzLabelPoints[left](P);
        \tkzLabelPoints[below left](O);
        \tkzLabelPoints[below right](F);
        \tkzLabelPoints[below](R);
    \end{tikzpicture}

    Agathe Deblouse affirme que $(PR)$ et $(TF)$ sont parallèles.
    Justifier qu'elle a raison.
\end{exercice*}
\begin{corrige}
    %\setcounter{partie}{0} % Pour s'assurer que le compteur de \partie est à zéro dans les corrigés
    \phantom{rrr}

    \begin{multicols}{2}
        Sur la figure ci-dessous :
        \begin{itemize}
            \item Les points $O$, $P$, $T$ sont alignés.
            \item Les points $O$, $R$, $F$ sont alignés.
        \end{itemize}
        \def\width{8}
        \def\hauteur{5}
        \begin{tikzpicture}[scale=0.8]        
            \draw[step=2mm, line width=0.1mm, black!30!white] (0,0) grid (\width,\hauteur);
            \draw[step=1cm, line width=0.3mm, black!50!white] (0,0) grid (\width,\hauteur);
            \coordinate (O) at (1,1);
            \coordinate (P) at (1,2);
            \coordinate (T) at (1,4);
            \coordinate (R) at (3,1);
            \coordinate (F) at (7,1);
            \tkzDrawSegment(O,T);
            \tkzDrawSegment(T,F);
            \tkzDrawSegment(F,O);
            \tkzDrawSegment(P,R);
            \tkzLabelPoints[above left](T);
            \tkzLabelPoints[left](P);
            \tkzLabelPoints[below left](O);
            \tkzLabelPoints[below right](F);
            \tkzLabelPoints[below](R);
        \end{tikzpicture}
    
        Agathe Deblouse affirme que $(PR)$ et $(TF)$ sont parallèles.    
        Justifier qu'elle a raison.
        \columnbreak

        % \smallskip
        {\color{red} 
        \Thales[Reciproque,Droites,Produit]{OTFPR}{1}{3}{2}{6}{}{}
        }
    \end{multicols}

\end{corrige}

